%%%% main.tex
%%%% Project specification paper based on the IJCAI--ECAI 26 template.

\typeout{Project paper using IJCAI--ECAI 26 formatting}

\documentclass{article}
\pdfpagewidth=8.5in
\pdfpageheight=11in

% Official IJCAI--ECAI 26 style file.
% Compile from this directory so the relative path resolves correctly.
\usepackage{../FormattingGuidelines-IJCAI-ECAI-26/ijcai26}

% Packages required/recommended by the official template.
\usepackage{times}
\usepackage{soul}
\usepackage{url}
\usepackage[hidelinks]{hyperref}
\usepackage[utf8]{inputenc}
\usepackage[small]{caption}
\usepackage{graphicx}
\usepackage{amsmath}
\usepackage{amsthm}
\usepackage{booktabs}
\usepackage{algorithm}
\usepackage{algorithmic}
\usepackage[switch]{lineno}

% Keep line numbers enabled for a submission. Disable for camera-ready.
\linenumbers

\urlstyle{same}

% Optional theorem environments.
\newtheorem{example}{Example}
\newtheorem{theorem}{Theorem}

% Required PDF metadata block. Keep unchanged.
\pdfinfo{
/TemplateVersion (IJCAI.2026.0)
}

\title{Projeto de IACH}

% For an anonymous submission, replace this block with the track-approved
% anonymous author declaration. For the camera-ready version, provide all
% author and affiliation information.
\author{
    Dinis Soveral Botelho Rosa
    \affiliations
    University of Coimbra / Department of Informatics Engineering
    \emails
    dinisrosa22@gmail.com
}

\begin{document}

\maketitle

\begin{abstract}
Briefly state the problem, motivation, proposed approach, and expected or
observed results. Keep the abstract self-contained and within the paper's
page limit.
\end{abstract}

\section{Introduction}
\label{sec:introduction}

Present the context and motivation of the project, define the problem being
addressed, and summarize the main contributions. End the section with a short
overview of the remaining paper.

\section{Problem Definition and Objectives}
\label{sec:problem}

Define the application context, stakeholders, inputs, outputs, constraints,
and success criteria. State the research questions or hypotheses that guide
the project.

\section{Background and Related Work}
\label{sec:related-work}

Introduce the concepts and previous approaches relevant to the project.
Discuss how the proposed work differs from or builds on existing methods.
Citations can be added with, for example, \verb|\cite{key}|.

\section{Proposed Approach}
\label{sec:approach}

Describe the system architecture, data flow, artificial intelligence methods,
and design decisions. Include mathematical notation, pseudocode, figures, or
tables when they clarify the proposal.

\subsection{System Architecture}
\label{sec:architecture}

Describe the main components and their interactions.

\subsection{Data and Preprocessing}
\label{sec:data}

Document the data sources, variables, preparation steps, and any ethical or
privacy considerations.

\subsection{Model or Reasoning Method}
\label{sec:method}

Specify the algorithm, model, prompts, rules, or reasoning procedure used by
the system, including relevant parameters and assumptions.

\section{Experimental Plan}
\label{sec:experiments}

Describe the experimental protocol, baselines, evaluation metrics, datasets,
and planned ablation or sensitivity analyses. Make the procedure reproducible.

\section{Expected Results and Discussion}
\label{sec:results}

Present results when available. Otherwise, state the expected outcomes and the
criteria that will be used to assess them. Discuss limitations and threats to
validity.

\section{Human-Centered and Ethical Considerations}
\label{sec:ethics}

Discuss human involvement, explainability, fairness, privacy, safety,
accessibility, and foreseeable misuse, as applicable to the project.

\section{Conclusion and Future Work}
\label{sec:conclusion}

Summarize the problem, approach, and main findings or expected contributions.
Identify the most relevant next steps.

% Optional in the final version; do not include for an anonymous submission
% when the selected track prohibits acknowledgements.
% \section*{Acknowledgements}
% Acknowledge support or contributions that are not part of the authorship.

% References may use the provided IJCAI named bibliography style.
\bibliographystyle{../FormattingGuidelines-IJCAI-ECAI-26/named}
\bibliography{references}

% Optional appendix. Remember that appendices count towards the applicable
% page limits unless the track instructions state otherwise.
% \appendix
% \section{Additional Material}
% Add supplementary details here.

\end{document}
